

\SourceLeaf{221}%
Since $a$, $b$, $c$ are to be regarded as vanishing in relation to this semicircle, and the points $A$, $B$ as therefore coinciding with $C$, this semicircle is the measure of the three angles of the triangle. Taken together, they therefore differ from $180^0$ by less than any given quantity.

It seems to me that, if one allows the concept of increase without end, this proof shows very simply that in every \textls{finite} rectilinear triangle the sum of the angles is $=180^0$. More precisely, it shows that the constant which, if Euclid's geometry were false, would have to be added to the sum of the angles to make it equal to $180^0$ is smaller than any given quantity. And since this can be proved for every triangle, the constant cannot depend on the size of the triangle either. \mbox{.\ .\ .\ .\ .\ .}\}

\bigskip
\hrule
\bigskip

\Needspace{7\baselineskip}
\begin{center}
[\textsc{Schumacher} \emph{to} \textsc{Gauss}. \emph{Altona, 29 June 1831.}]
\end{center}

\{\mbox{.\ .\ .\ .\ .\ .}\ .\ .\ There was just one thing I missed in your letter, namely your judgment on my proof that the sum of the angles of a rectilinear triangle differs from $180^0$ only by a quantity smaller than any given quantity. You can readily imagine how important your judgment is to me, since you so readily detect every weakness in a proof. I have told no one about it besides you, my assistant, and Professor \textsc{Hansen} of Seeberg. None of us can detect a paralogism.

Suppose someone thinks it necessary to prove that the vertices of a finite triangle may be regarded as coincident centers of a circle of infinite radius (called \textls{infinite} brevitatis causa). Although I do not think this necessary, the proof can be given rigorously.

It seems to me that if two points are a finite distance apart, this distance must be set $=0$ in relation to an infinite line. The points therefore coincide when considered in relation to this infinite line. \mbox{.\ .\ .\ .\ .\ .}\ .\}

\bigskip
\hrule
\bigskip
